\casechapter{Axiomatizing Chance}{Kolmogorov and the hidden parameter}{Axiomatization}{ch:kolmogorov}
\begin{refsection}[chapters/03-kolmogorov.bib]
\newcommand{\kolP}{\mathbb{P}}
\newcommand{\kolsymdiff}{\mathbin{\triangle}}

\begin{chaptersummary}
Before 1933 probability had paradoxes that nobody could settle, because nobody could say what a question had failed to specify. A random chord of a circle is long with probability $\frac13$, $\frac12$, or $\frac14$; Borel's strong law spoke of infinitely many coin tosses that no theory defined; a uniform point on a sphere, told that it lies on a great circle, is uniform on that circle or not, depending on how one computes. Kolmogorov made a probability a measure of total mass one, a random variable a measurable function, expectation an integral, independence a product rule, and conditioning a projection onto a $\sigma$-algebra. Each paradox turned out to hide an undeclared parameter, a measure, a probability space, or a $\sigma$-algebra, and the axioms force it into the open. After the move Borel's normal numbers follow from one strong law in a few lines, the zero--one law takes half a page, and the sphere paradox becomes two correct answers to two different questions.
\end{chaptersummary}

\section{Before: paradoxes without parameters}\label{sec:kolmogorov:before}

For two centuries probability was a calculus of equally possible cases. Its great theorems, from Bernoulli's law of large numbers to the normal approximations of de~Moivre and Laplace, concerned finitely many trials, and its basic notion stayed informal. Around 1900 three kinds of questions pushed the calculus past its definitions: questions about a point or a line chosen at random from a continuum, questions about infinitely many trials, and questions that condition on an event of probability zero. In each, one reasonable computation gave one answer and another gave a different one. In 1900 Hilbert's sixth problem asked that the physical sciences in which mathematics plays an important part be treated by axioms, and it put probability, with mechanics, in the first rank \parencite{Hilbert1902}.\sidenote{Hilbert added that the logical study of the axioms of probability should go with a rigorous development of the method of mean values in mathematical physics, in particular in the kinetic theory of gases. His footnote points to Bohlmann's work on insurance mathematics.}

\subsection{Random chords}\label{sec:kolmogorov:chords}

In its usual form, the paradox of \textcite{Bertrand1889} asks for the probability that a chord drawn at random in a circle is longer than the side of the inscribed equilateral triangle.\index{Bertrand's paradox} Take the unit circle. The side has length $\sqrt3$, and a chord is longer than $\sqrt3$ exactly when its distance from the center is less than $\frac12$. Three ways of drawing a chord at random give three answers (\cref{fig:kolmogorov:bertrand}).
\begin{enumerate}
\item Random endpoints. Fix one endpoint, by symmetry, and choose the other uniformly on the circle. The chord is long when the second endpoint lies on the arc of angle $\frac{2\pi}3$ opposite the first: probability $\frac13$.
\item Random radius. Fix the direction of the perpendicular from the center to the chord, by symmetry, and choose the chord's distance from the center uniformly in $[0,1]$: probability $\frac12$.
\item Random midpoint. Choose the chord's midpoint uniformly in the disk. The chord is long when the midpoint lies in the concentric disk of radius $\frac12$: probability $\frac14$.
\end{enumerate}
Each computation is correct. Bertrand's verdict, as \textcite{Jaynes1973} reports it, was that none of the three answers is false and none is exact: the question is ill posed. That is a diagnosis without a cure. Nothing in the calculus of equally possible cases said what a well-posed question about a continuum must supply.

\begin{figure}[ht]
\centering
\begin{tikzpicture}[x=1.05cm,y=1.05cm,baseline=0pt]
  \draw[sandline] (270:1) -- (30:1) -- (150:1) -- cycle;
  \draw[muted] (0,0) circle (1);
  \draw[seadeep,line width=2pt] (30:1) arc (30:150:1);
  \draw[sea,line width=0.9pt] (270:1) -- (105:1);
  \draw[sea,line width=0.9pt] (270:1) -- (195:1);
  \fill[seadeep] (270:1) circle (1.6pt);
  \node[font=\small] at (0,-1.45) {(i)\enspace$\frac13$};
\end{tikzpicture}\hspace{0.9cm}%
\begin{tikzpicture}[x=1.05cm,y=1.05cm,baseline=0pt]
  \draw[sandline] (270:1) -- (30:1) -- (150:1) -- cycle;
  \draw[muted] (0,0) circle (1);
  \draw[muted] (0,0) -- (0,1);
  \draw[seadeep,line width=2pt] (0,0) -- (0,0.5);
  \draw[sea,line width=0.9pt] (-0.975,0.22) -- (0.975,0.22);
  \fill[seadeep] (0,0.22) circle (1.4pt);
  \node[font=\small] at (0,-1.45) {(ii)\enspace$\frac12$};
\end{tikzpicture}\hspace{0.9cm}%
\begin{tikzpicture}[x=1.05cm,y=1.05cm,baseline=0pt]
  \fill[seamist] (0,0) circle (0.5);
  \draw[sandline] (270:1) -- (30:1) -- (150:1) -- cycle;
  \draw[muted] (0,0) circle (1);
  \draw[sea,line width=0.9pt] (-0.342,0.94) -- (0.842,-0.54);
  \fill[seadeep] (0.25,0.2) circle (1.4pt);
  \node[font=\small] at (0,-1.45) {(iii)\enspace$\frac14$};
\end{tikzpicture}
\caption{Bertrand's three chords. (i) From a fixed endpoint, the chord is long when the other endpoint falls on the thick arc. (ii) It is long when its distance from the center falls on the thick half of the radius. (iii) It is long when its midpoint falls in the shaded disk of radius $\frac12$.}
\label{fig:kolmogorov:bertrand}
\end{figure}

\subsection{Infinitely many coins}\label{sec:kolmogorov:coins}

\textcite{Borel1909} asserted that almost every number is \term{absolutely normal}: for almost every $x\in[0,1]$, in every base, every block of digits occurs in the expansion of $x$ with its fair limiting frequency.\index{normal number} In particular, the proportion of ones among the first $n$ binary digits of $x$ tends to $\frac12$. He reasoned with ``denumerable probabilities'', extending the rules of total and compound probability to countably many trials, and he deliberately set Lebesgue measure aside.\sidenote{Borel remarked that measure theory would give a proof, but declined to rely on it, for philosophical reasons that concerned countable additivity \parencite{ShaferVovk2006}.} His proof had gaps; rigorous proofs were given by Faber in 1910 and by Hausdorff in 1914 \parencite{ShaferVovk2006}. The obstacle was not the estimates but the objects. An event such as ``the frequency of heads tends to $\frac12$'' depends on infinitely many tosses, while the classical calculus assigns probabilities only to events about finitely many. The aside shows what a rigorous proof looked like without a theory of probability: a computation with lengths of intervals.

\begin{aside}[Borel's theorem in base two, by counting and covering]
For $x\in[0,1)$ let $d_k(x)$ be its $k$-th binary digit, taking the expansion that does not end in ones, and let $r_k=2d_k-1$ and $S_n=r_1+\dots+r_n$. We show that the set of $x$ with $S_n(x)/n\not\to0$, which is the set where the frequency of ones does not tend to $\frac12$, can be covered, for every $\eps>0$, by countably many intervals of total length less than $\eps$.

\emph{Counting.} The function $S_n$ is constant on each of the $2^n$ dyadic intervals $[j2^{-n},(j+1)2^{-n})$, and as $x$ runs over them, $(r_1,\dots,r_n)$ runs over the $2^n$ sign patterns once each. Expand $S_n^4=\sum_{i,j,k,l}r_ir_jr_kr_l$ and sum over the patterns: a term survives exactly when each index appears an even number of times, which leaves the $n$ terms $r_i^4$ and the $3n(n-1)$ terms $r_i^2r_j^2$ with $i\ne j$, each summing to $2^n$. So the dyadic intervals of rank $n$ on which $\abs{S_n}\ge t$ have total length at most $2^{-n}t^{-4}\cdot2^n(3n^2-2n)\le3n^2t^{-4}$.

\emph{Covering.} Let $B_n$ be the union of the dyadic intervals of rank $n$ on which $\abs{S_n}\ge n^{7/8}$; its total length is at most $3n^{-3/2}$. If $x$ lies in only finitely many $B_n$, then $\abs{S_n(x)}/n<n^{-1/8}$ for large $n$, and $S_n(x)/n\to0$. Every other $x$ lies in $\bigcup_{n\ge N}B_n$ for each $N$, a countable union of intervals of total length at most $3\sum_{n\ge N}n^{-3/2}$, which tends to $0$.

\emph{Where the difficulty lies.} The argument is a computation on one model, the digits of a point of $[0,1)$, with the counting done by hand. Change the model, to digits in base three, to blocks of digits rather than single digits, or to a biased coin, and the counting must be redone; for overlapping blocks, whose indicators are not independent, it becomes awkward. The probabilistic content, independence and fourth moments and summability, is real, but it has nowhere to live.
\end{aside}

\textcite{Steinhaus1923} made the identification systematic: Borel's denumerable probabilities for coin tossing can be realized by Lebesgue measure on $[0,1]$, the tosses being the binary digits. That settles coin tossing. But every new model, a die, a biased coin, a continuous measurement, infinitely many of them at once, needs its own encoding into the unit interval, and every theorem must be proved again inside it.

\subsection{Conditioning on what did not happen}\label{sec:kolmogorov:conditioning}

Conditional probability was defined as a ratio, $\kolP(A\mid B)=\kolP(A\cap B)/\kolP(B)$, which needs $\kolP(B)>0$. For quantities with densities the textbooks used a ratio of densities: if $(X,Y)$ has the joint density $f$ and $Y$ has the density $f_Y(y)=\int f(x,y)\dd x$, then the conditional density of $X$ given $Y=y$ is $f(x,y)/f_Y(y)$. The familiar rules can then be checked by calculus, under hypotheses that keep the ratios meaningful. Here is the rule for iterated conditioning.

\begin{proposition}[The tower rule with densities]\label[proposition]{prop:kolmogorov:tower}
Let $(X,Y,Z)$ take values in $[0,1]^3$ and have a continuous joint density $f$ that is positive on $[0,1]^3$, and use the continuous marginal densities $f_{XY}(x,y)=\int_0^1f\dd z$, $f_{YZ}(y,z)=\int_0^1f\dd x$, and $f_Y(y)=\int_0^1\!\int_0^1f\dd x\dd z$. Put
\[
  e(y,z)=\frac{\int_0^1xf(x,y,z)\dd x}{f_{YZ}(y,z)}\quad\text{and}\quad e(y)=\frac{\int_0^1xf_{XY}(x,y)\dd x}{f_Y(y)} .
\]
Then $\int_0^1e(y,z)\,\frac{f_{YZ}(y,z)}{f_Y(y)}\dd z=e(y)$ for every $y\in[0,1]$.
\end{proposition}

\begin{proof}
Since $f$ is continuous and positive on the compact cube, these marginal densities are continuous and positive, and every integral below is an integral of a continuous function. Clearing the denominator of $e(y,z)$ and integrating over $z$ gives $\int_0^1e(y,z)f_{YZ}(y,z)\dd z=\int_0^1\!\int_0^1xf(x,y,z)\dd x\dd z$. Exchanging the order of integration, which is allowed for continuous functions on a square, turns the right side into $\int_0^1xf_{XY}(x,y)\dd x=e(y)f_Y(y)$. Divide by $f_Y(y)>0$.
\end{proof}

The proof is short, but the hypotheses carry the weight: a joint density, continuous and positive, on a compact cube. Discrete or mixed variables need a separate calculation. And the formula invites us to read $f(x,y)/f_Y(y)$ as the law of $X$ given the null event $\{Y=y\}$ alone. The next example shows that this reading cannot be right: the same kind of null event gets different laws depending on the variable that defines it.

Let $U$ be a uniformly distributed point on the unit sphere, with longitude $\Lambda\in[-\pi,\pi)$ and latitude $\Phi\in[-\frac\pi2,\frac\pi2]$. The area element is $\cos\varphi\dd\varphi\dd\lambda$, so $(\Lambda,\Phi)$ has the joint density $\frac{\cos\varphi}{4\pi}$, and $\Phi$ and $\Lambda$ have the densities $\frac12\cos\varphi$ and $\frac1{2\pi}$. Condition on the equator, $\Phi=0$: the ratio of densities gives $\Lambda$ the conditional density $\frac{1}{4\pi}\big/\frac12=\frac1{2\pi}$, the uniform law. Condition on a meridian, $\Lambda=\lambda_0$: it gives $\Phi$ the conditional density $\frac{\cos\varphi}{4\pi}\big/\frac1{2\pi}=\frac12\cos\varphi$, which favors the equator over the poles. Yet a rotation of the sphere, which preserves the uniform distribution, carries the meridian onto half of the equator. Two half great circles, congruent by a symmetry of the problem, receive different conditional laws.\index{Borel--Kolmogorov paradox} A version of this paradox appears in Bertrand's book, Borel discussed it in 1909, and Kolmogorov explained it.\sidenote{Bertrand's version concerns two random points on a sphere; Borel's discussion uses thin bundles of great circles \parencite{ShaferVovk2006}.}

In each example the computation is sound and the question is not. Something that the computation needs and the question does not say has been chosen silently: a measure on the set of chords, a probability for events about infinite sequences of tosses, or the family of conditions of which ``$\Phi=0$'' or ``$\Lambda=\lambda_0$'' is one member. Call it the \term{hidden parameter}.\index{hidden parameter}

\section{The move: probability is a measure}\label{sec:kolmogorov:move}

Kolmogorov's \emph{Grundbegriffe der Wahrscheinlichkeitsrechnung} \parencite{Kolmogorov1933}, a monograph of sixty-two pages,\sidenote{It appeared in Springer's series \emph{Ergebnisse der Mathematik und ihrer Grenzgebiete}; a Russian translation followed in 1936, and English ones in 1950 and 1956.} takes the measure theory of \cref{ch:lebesgue}, by then extended from the line to abstract sets, and declares that a probability is a measure of total mass one. Most of the definitions were not new; the move was to make measure the whole foundation, and to add a new definition of conditioning.\sidenote{Kolmogorov stated the axioms on a field of sets, with finite additivity and a continuity axiom ($\kolP(A_n)\to0$ whenever $A_1\supseteq A_2\supseteq\cdots$ have empty intersection); together these give countable additivity, and Carathéodory's theorem extends $\kolP$ to the generated $\sigma$-algebra.}

\begin{definition}[Probability spaces and random variables]\label[definition]{def:kolmogorov:space}
A \term{probability space} is a triple $(\Omega,\cF,\kolP)$: a set $\Omega$, a $\sigma$-algebra $\cF$ of subsets of $\Omega$, called \term{events}, and a measure $\kolP$ on $\cF$ with $\kolP(\Omega)=1$. A \term{random variable} is a measurable function $X\colon\Omega\to\R$; its \term{distribution} is the measure $B\mapsto\kolP(X\in B)$ on the Borel sets of $\R$, and its \term{expectation} is $\E X=\int X\dd\kolP$ whenever $X$ is integrable or nonnegative; we use the same notation for measurable functions with values in $[0,\infty]$. A property that fails only on a subset of an event of probability zero holds \term{almost surely} (a.s.).
\end{definition}
\index{probability space}

\begin{definition}[Independence]\label[definition]{def:kolmogorov:independence}
The $\sigma$-algebras $\cF_1,\cF_2,\ldots\subseteq\cF$ are \term{independent} if $\kolP(A_1\cap\dots\cap A_k)=\kolP(A_1)\cdots\kolP(A_k)$ whenever $A_j\in\cF_{i_j}$ for distinct indices $i_1,\dots,i_k$. Random variables $X_1,X_2,\dots$ are independent if the $\sigma$-algebras $\sigma(X_i)=\{X_i^{-1}(B):B\text{ Borel}\}$ are, and events are independent if their indicator functions are.
\end{definition}
\index{independence}

We use two standard facts about independence freely \parencite{Williams1991,Durrett2019}. If $X$ and $Y$ are independent and integrable, then $XY$ is integrable and $\E[XY]=\E X\,\E Y$. And independence survives grouping: if $X_1,X_2,\dots$ are independent and $I_1,I_2,\dots$ are disjoint sets of indices, then the $\sigma$-algebras generated by the groups $\{X_i:i\in I_j\}$ are independent, so random variables that are measurable with respect to these $\sigma$-algebras are independent.

A definition of this kind is worth something only when its objects exist. For infinite families, existence is a theorem, the ``fundamental theorem'' of the \emph{Grundbegriffe}; it is the one major theorem that this chapter uses without proof.\index{Kolmogorov's extension theorem}

\begin{theorem}[Kolmogorov's extension theorem]\label{thm:kolmogorov:extension}
Let $T$ be a set, and for each finite $S\subseteq T$ let $\mu_S$ be a probability measure on the Borel sets of $\R^S$, consistently: when $S\subseteq S'$, $\mu_S$ is the image of $\mu_{S'}$ under the coordinate projection $\R^{S'}\to\R^S$. Then there is a unique probability measure on $\R^T$, with the $\sigma$-algebra generated by the coordinate functions, whose image on each $\R^S$ is $\mu_S$. In particular, for all probability measures $\mu_1,\mu_2,\dots$ on $\R$ there is a probability space carrying independent random variables $X_1,X_2,\dots$ with distributions $\mu_1,\mu_2,\dots$.
\end{theorem}

For coin tossing the theorem is not needed, because \cref{ch:lebesgue} already built the space.

\begin{proposition}[Digits are independent]\label[proposition]{prop:kolmogorov:digits}
Let $b\ge2$. Under Lebesgue measure on $[0,1)$, the base-$b$ digits $d_1,d_2,\dots$ (taking expansions that do not end in the digit $b-1$) are independent, and each is uniformly distributed on $\{0,\dots,b-1\}$.
\end{proposition}

\begin{proof}
For digits $e_1,\dots,e_n$, the set $\{d_1=e_1,\dots,d_n=e_n\}$ is an interval of length $b^{-n}$. So the first $n$ digits are uniformly distributed on $\{0,\dots,b-1\}^n$, and summing over the digits not under consideration gives $\kolP(d_{i_1}=e_1,\dots,d_{i_k}=e_k)=b^{-k}$ for distinct indices. An event of $\sigma(d_i)$ has the form $\{d_i\in E\}$ with $E\subseteq\{0,\dots,b-1\}$, a disjoint union of events $\{d_i=e\}$, so the product rule extends to these events by additivity.
\end{proof}

Conditioning is the last notion, and the subtle one.

\begin{definition}[Conditional expectation]\label[definition]{def:kolmogorov:condexp}
Let $\cG\subseteq\cF$ be a $\sigma$-algebra and $X$ an integrable random variable. A \term{conditional expectation} of $X$ given $\cG$ is a $\cG$-measurable integrable random variable $Y$ with $\E[Y\ind_G]=\E[X\ind_G]$ for every $G\in\cG$. We write $Y=\E[X\mid\cG]$, $\kolP(A\mid\cG)=\E[\ind_A\mid\cG]$, and $\E[X\mid Z]=\E[X\mid\sigma(Z)]$ for a random variable $Z$.
\end{definition}
\index{conditional expectation}

\textcite{Kolmogorov1933} proved that conditional expectations exist with the Radon--Nikodym theorem of \textcite{Nikodym1930};\sidenote{Radon had proved the theorem for measures on $\R^n$ in 1913; Nikodym's version holds on abstract measure spaces.} we shall prove it for square-integrable $X$ by projection. A conditional expectation is a random variable, determined only up to events of probability zero. That last clause is where the sphere paradox will dissolve.

With these definitions, each paradox becomes a question with a missing argument. ``A random chord'' needs a probability measure on the set of chords. ``Infinitely many tosses'' needs a probability space that carries them, such as Lebesgue measure on $[0,1)$ or the measure of \cref{thm:kolmogorov:extension}. And ``given that $U$ lies on the great circle $C$'' needs a $\sigma$-algebra in which the information ``$U\in C$'' is one possible observation among others. Kolmogorov put the third point sharply: a conditional probability with respect to an isolated hypothesis of probability zero is inadmissible; a distribution on a meridian makes sense only when the meridian is regarded as an element of the partition of the sphere into meridians \parencite[Chapter~V, §2]{Kolmogorov1956}.

The new object is humble: a $\sigma$-algebra, the list of questions whose answers we are told. A partition of the sphere into circles of latitude is one experiment, a partition into meridians is another, and an observation is a cell of the experiment that produced it. Kolmogorov's move was to stop conditioning on events and to condition on experiments.

What would one have had to notice? That the rigorous versions of Borel's arguments, by Faber, Hausdorff, and Steinhaus, were already measure theory, that expectation was already Lebesgue's integral, and that the only notion probability adds to measure is independence, together with conditioning, which measure theory can also express. The axioms say nothing new about measures. They say that nothing else is needed.

\section{After: one strong law for every model}\label{sec:kolmogorov:limits}

The first dividend is that a limit theorem, once proved, applies to every model at once.\index{Borel--Cantelli lemma}

\begin{lemma}[Borel--Cantelli]\label[lemma]{lem:kolmogorov:bc}
If $A_1,A_2,\ldots$ are events with $\sum_n\kolP(A_n)<\infty$, then almost surely only finitely many of them occur.
\end{lemma}

\begin{proof}
The number $N=\sum_n\ind_{A_n}$ of events that occur has $\E N=\sum_n\kolP(A_n)<\infty$, by term-by-term integration (\cref{cor:lebesgue:series}), so $N$ is finite almost surely.
\end{proof}

Borel used the lemma for independent events; \textcite{Cantelli1917} applied the convergent half, which needs no independence, in his own strong law. The proof of the strong law below is the standard one by fourth moments \parencite{Williams1991}.\sidenote{The theorem is a special case of Kolmogorov's strong law of 1930, which needs only $\sum_n\Var X_n/n^2<\infty$.}\index{strong law of large numbers}

\begin{theorem}[Strong law of large numbers, fourth-moment form]\label{thm:kolmogorov:slln}
Let $X_1,X_2,\dots$ be independent random variables, and let $\mu$ and $K$ be constants with $\E X_i=\mu$ and $\E(X_i-\mu)^4\le K$ for every $i$. Then $S_n=X_1+\dots+X_n$ satisfies $S_n/n\to\mu$ almost surely.
\end{theorem}

\begin{proof}
The plan is to bound the fourth moment of the $n$-th sum by $3Kn^2$, which makes $\sum_n\E(S_n/n)^4$ finite. Replacing $X_i$ by $X_i-\mu$ and $S_n$ by $S_n-n\mu$, we may assume $\mu=0$. Expand $\E S_n^4=\sum_{i,j,k,l}\E[X_iX_jX_kX_l]$; the products are integrable by Hölder's inequality. If some index appears exactly once, then the expectation factors, by independence and grouping, and contains a factor $\E X_i=0$. Two kinds of terms survive: the $n$ terms $\E X_i^4\le K$, and the $3n(n-1)$ terms $\E[X_i^2X_j^2]=\E X_i^2\,\E X_j^2$ with $i\ne j$, each at most $K$ because $(\E X_i^2)^2\le\E X_i^4$. Hence
\[
  \E S_n^4\le nK+3n(n-1)K\le3Kn^2 .
\]
By term-by-term integration, $\E\sum_n(S_n/n)^4\le\sum_n3K/n^2<\infty$. So $\sum_n(S_n/n)^4$ is finite almost surely, and its terms tend to $0$.
\end{proof}

\begin{corollary}[Borel's normal number theorem]\label[corollary]{cor:kolmogorov:normal}
Almost every $x\in[0,1)$ is absolutely normal: for every base $b\ge2$ and every block $w$ of $\ell$ digits in base $b$, the proportion of positions $m\le n$ at which $w$ begins in the base-$b$ expansion of $x$ tends to $b^{-\ell}$ as $n\to\infty$.
\end{corollary}

\begin{proof}
Fix $b$ and $w$. By \cref{prop:kolmogorov:digits} the digits are independent and uniform. Let $I_m$ be the indicator that $w$ begins at position $m$; it is a function of the digits in positions $m,\dots,m+\ell-1$, and $\E I_m=b^{-\ell}$. For each residue $r$ modulo $\ell$, the $I_m$ with $m\equiv r$ depend on disjoint blocks of digits, so they are independent, by grouping, and bounded, and \cref{thm:kolmogorov:slln} says that their averages tend to $b^{-\ell}$ almost surely. The average of all $I_m$ with $m\le n$ is a weighted average of these $\ell$ averages, with weights tending to $\frac1\ell$, so it also tends to $b^{-\ell}$. There are countably many pairs $(b,w)$, and a countable union of null sets is null.
\end{proof}

Compare the aside. The counting and covering have been replaced by one theorem, proved once, and the corollary covers every base and every block, including the overlapping blocks that made the counting awkward. The reason is not cleverness. The strong law is now a theorem about every independent sequence with a common mean and bounded fourth moments, in every probability space, so a new model costs nothing.

The second dividend is a theorem that the old setting could not state in general. It needs one lemma, which says that events in a generated $\sigma$-algebra are approximated by events in the generating algebra.

\begin{lemma}[Approximation]\label[lemma]{lem:kolmogorov:approx}
Let $\cA_0$ be a family of events that contains $\Omega$ and is closed under complements and finite unions. Then for every $A\in\sigma(\cA_0)$ and every $\eps>0$ there is $B\in\cA_0$ with $\kolP(A\kolsymdiff B)<\eps$.
\end{lemma}

\begin{proof}
Let $\cD$ be the family of events $A$ with this property. It contains $\cA_0$, and it is closed under complements because $A^c\kolsymdiff B^c=A\kolsymdiff B$. It is closed under countable unions. Indeed, let $A=\bigcup_iA_i$ with $A_i\in\cD$, and choose $N$ with $\kolP(A\setminus\bigcup_{i\le N}A_i)<\frac\eps2$, by continuity from below (\cref{lem:lebesgue:continuity}), and $B_i\in\cA_0$ with $\kolP(A_i\kolsymdiff B_i)<\eps2^{-i-1}$. The set $B=\bigcup_{i\le N}B_i$ lies in $\cA_0$, and $A\kolsymdiff B$ is contained in $(A\setminus\bigcup_{i\le N}A_i)\cup\bigcup_{i\le N}(A_i\kolsymdiff B_i)$, which has probability less than $\eps$. So $\cD$ is a $\sigma$-algebra that contains $\cA_0$, and therefore $\cD\supseteq\sigma(\cA_0)$.
\end{proof}

\begin{theorem}[Kolmogorov's zero--one law]\label{thm:kolmogorov:zeroone}
Let $X_1,X_2,\dots$ be independent, and let $\cT=\bigcap_n\sigma(X_n,X_{n+1},\dots)$ be their \term{tail $\sigma$-algebra}. Then every $A\in\cT$ has $\kolP(A)\in\{0,1\}$.
\end{theorem}
\index{zero--one law}

\begin{proof}
The plan is to show that $A$ is independent of itself. The family $\cA_0=\bigcup_n\sigma(X_1,\dots,X_n)$ contains $\Omega$, is closed under complements and finite unions, and generates $\sigma(X_1,X_2,\dots)$, which contains $A$. By \cref{lem:kolmogorov:approx} there are $B_k\in\cA_0$ with $\kolP(A\kolsymdiff B_k)\to0$. Each $B_k$ lies in some $\sigma(X_1,\dots,X_{n_k})$, while $A$ lies in $\sigma(X_{n_k+1},X_{n_k+2},\dots)$, so $A$ and $B_k$ are independent by grouping, and $\kolP(A\cap B_k)=\kolP(A)\kolP(B_k)$. As $k\to\infty$, both $\kolP(A\cap B_k)$ and $\kolP(B_k)$ tend to $\kolP(A)$, since each differs from it by at most $\kolP(A\kolsymdiff B_k)$. Hence $\kolP(A)=\kolP(A)^2$.
\end{proof}

For each $m$, the averages $\frac1nS_n$ and $\frac1n(X_m+\dots+X_n)$ differ by a term that tends to $0$, so the convergence and the $\limsup$ of $\frac1nS_n$ are determined by $X_m,X_{m+1},\dots$ for every $m$; likewise $\sum_nX_n$ converges exactly when $\sum_{n\ge m}X_n$ does. Before a single estimate, the zero--one law says that $\frac1nS_n$ converges with probability $0$ or $1$, and that $\limsup_n\frac1nS_n$ is almost surely equal to a constant in $[-\infty,\infty]$. Under the hypotheses of \cref{thm:kolmogorov:slln}, the strong law says which.

\section{After: conditioning is a projection}\label{sec:kolmogorov:projection}

The third dividend is conditioning without densities. The space $L^2(\Omega,\cG,\kolP)$ of square-integrable $\cG$-measurable random variables sits inside $L^2(\Omega,\cF,\kolP)$, and conditioning is the orthogonal projection onto it.\sidenote{Completeness is what \cref{ch:lebesgue} bought; in an incomplete space of random variables the projection could fail to exist.} We use the projection theorem for Hilbert spaces: if $V$ is a closed subspace of a Hilbert space $H$, then each $x\in H$ has a unique $y\in V$ with $x-y$ orthogonal to $V$ \parencite{Williams1991}.\index{tower property}

\begin{theorem}[Conditional expectation as a projection]\label{thm:kolmogorov:projection}
Let $\cG\subseteq\cF$ be a $\sigma$-algebra and $X\in L^2(\Omega,\cF,\kolP)$. The orthogonal projection of $X$ onto $L^2(\Omega,\cG,\kolP)$ is a conditional expectation of $X$ given $\cG$, and two conditional expectations of an integrable random variable given $\cG$ agree almost surely. Moreover, for $X\in L^2$:
\begin{enumerate}
\item (tower property) if $\cH\subseteq\cG$ is a $\sigma$-algebra, then $\E\bigl[\E[X\mid\cG]\bigm|\cH\bigr]=\E[X\mid\cH]$;
\item the expectation is preserved: $\E\bigl[\E[X\mid\cG]\bigr]=\E X$;
\item if $Z$ is bounded and $\cG$-measurable, then $\E[ZX\mid\cG]=Z\,\E[X\mid\cG]$.
\end{enumerate}
\end{theorem}

\begin{proof}
The plan is to let completeness supply the projection and orthogonality supply the defining identity. By the Riesz--Fischer theorem (\cref{thm:lebesgue:rieszfischer}) applied to $(\Omega,\cG,\kolP)$, the space $L^2(\cG)$ is complete, hence closed in $L^2(\cF)$, so the projection of $X$ onto it exists. Choose a $\cG$-measurable representative $Y$; it is integrable, because $\E\abs{Y}\le(\E Y^2)^{1/2}$, and $X-Y$ is orthogonal to $L^2(\cG)$; since $\ind_G\in L^2(\cG)$ for $G\in\cG$, this gives $\E[(X-Y)\ind_G]=0$, the defining identity.

For uniqueness, let $Y$ and $Y'$ be conditional expectations of an integrable $X$. Then $W=Y-Y'$ is $\cG$-measurable and $\E[W\ind_G]=0$ for every $G\in\cG$. The choice $G=\{W>0\}$ gives $\E W^+=0$, and the choice $G=\{W<0\}$ gives $\E W^-=0$, so $W=0$ almost surely.

Part (i) holds because $L^2(\cH)\subseteq L^2(\cG)$: the difference $X-\E[X\mid\cG]$ is orthogonal to $L^2(\cG)$, hence to $L^2(\cH)$, so $X$ and $\E[X\mid\cG]$ have the same projection onto $L^2(\cH)$. Part (ii) is the defining identity with $G=\Omega$. For (iii), the random variable $Z\,\E[X\mid\cG]$ is $\cG$-measurable and square integrable, and for $G\in\cG$ the function $Z\ind_G$ lies in $L^2(\cG)$, so orthogonality gives $\E[ZX\ind_G]=\E\bigl[\E[X\mid\cG]\,Z\ind_G\bigr]$.
\end{proof}

Existence extends to every integrable $X$. For $X\ge0$, conditional expectations $Y_n$ of $\min(X,n)$ are nonnegative and increase with $n$ almost surely, because conditional expectation preserves order (Exercise~\ref{ex:kolmogorov:condexp}); so $Y=\limsup_nY_n$, set to $0$ where it is infinite, is $\cG$-measurable, and monotone convergence gives $\E[Y\ind_G]=\E[X\ind_G]$ for $G\in\cG$, with $\E Y=\E X<\infty$. For general $X$, use the positive and negative parts. Kolmogorov obtained existence in one stroke from the Radon--Nikodym theorem.

Compare \cref{prop:kolmogorov:tower}. The tower property needs only the definition: for integrable $X$ and $H\in\cH\subseteq\cG$, both $\E[X\mid\cH]$ and $\E\bigl[\E[X\mid\cG]\bigm|\cH\bigr]$ are $\cH$-measurable with the same integral as $X$ over $H$, so they agree almost surely. That is one line, with no densities, no continuity, no positivity, and no compact cube, and it holds for every pair of nested $\sigma$-algebras, whatever the random variables are. The same theorem dissolves the sphere paradox.

\begin{proposition}[Two conditionings on the sphere]\label[proposition]{prop:kolmogorov:sphere}
Let $U$ be uniformly distributed on the unit sphere, with longitude $\Lambda$ and latitude $\Phi$, the poles being ignored as a null set. Then $\Lambda$ and $\Phi$ are independent, $\Lambda$ is uniform on $[-\pi,\pi)$, and $\Phi$ has density $\frac12\cos\varphi$ on $[-\frac\pi2,\frac\pi2]$. Consequently, for every bounded measurable $h$, almost surely,
\begin{gather*}
  \E[h(\Lambda,\Phi)\mid\Phi]=\frac1{2\pi}\int_{-\pi}^{\pi}h(\lambda,\Phi)\dd\lambda,\\
  \E[h(\Lambda,\Phi)\mid\Lambda]=\frac12\int_{-\pi/2}^{\pi/2}h(\Lambda,\varphi)\cos\varphi\dd\varphi .
\end{gather*}
\end{proposition}

\begin{proof}
The area element in these coordinates is $\cos\varphi\dd\varphi\dd\lambda$ and the total area is $4\pi$, so $(\Lambda,\Phi)$ has the joint density $\frac1{2\pi}\cdot\frac12\cos\varphi$, a product of the two stated densities; this gives independence and the distributions. The right side of the first formula is $g(\Phi)$, with $g(\varphi)=\frac1{2\pi}\int h(\lambda,\varphi)\dd\lambda$, a Borel function by Fubini's theorem; so it is $\sigma(\Phi)$-measurable, and for a Borel set $B$, Fubini's theorem gives
\[
  \E\bigl[h(\Lambda,\Phi)\ind\{\Phi\in B\}\bigr]=\int_Bg(\varphi)\,\tfrac12\cos\varphi\dd\varphi=\E\bigl[g(\Phi)\ind\{\Phi\in B\}\bigr].
\]
The second formula follows in the same way, with the roles of $\Lambda$ and $\Phi$ exchanged.
\end{proof}

Both formulas are correct, and they do not contradict each other. The first, read on the null event $\{\Phi=0\}$, makes the longitude uniform on the equator; the second, read on the null event $\{\Lambda=\lambda_0\}$, gives the latitude the density $\frac12\cos\varphi$ on a meridian. They are values of two conditional expectations, given two $\sigma$-algebras, on two null sets, and a conditional expectation is determined only up to null sets: its values on one null set mean something only as part of the whole family. The question ``what is the law of $U$ given that $U$ lies on the great circle $C$?'' has no answer, because $\{U\in C\}$ has probability zero. The question ``what is the law of $U$ given its circle of latitude?'' has one, so does ``given its meridian'', and the two answers differ on two half great circles that a rotation exchanges. The parameter that the paradox hid is the $\sigma$-algebra.

The answers have a geometric meaning (\cref{fig:kolmogorov:sphere}). Conditioning on the band $\{\abs\Phi<\eps\}$, an event of positive probability, makes the longitude exactly uniform, for every $\eps$; conditioning on the lune $\{\abs{\Lambda-\lambda_0}<\eps\}$, the difference taken modulo $2\pi$, gives the latitude exactly the density $\frac12\cos\varphi$, for every $\eps$, because the lune is widest at the equator. Shrinking the band and the lune does not remove the difference; it is built into their shapes (Exercise~\ref{ex:kolmogorov:lunes}).

\begin{figure}[ht]
\centering
\begin{tikzpicture}[x=1.4cm,y=1.4cm,baseline=0pt]
  \def\al{18}\def\ep{14}
  \fill[seamist] plot[domain=-90:90,samples=60,smooth] ({cos(\ep)*sin(\x)},{cos(\al)*sin(\ep)-sin(\al)*cos(\ep)*cos(\x)}) -- plot[domain=90:-90,samples=60,smooth] ({cos(\ep)*sin(\x)},{-cos(\al)*sin(\ep)-sin(\al)*cos(\ep)*cos(\x)}) -- cycle;
  \draw[sea!60,line width=0.4pt] plot[domain=-90:90,samples=60,smooth] ({cos(\ep)*sin(\x)},{cos(\al)*sin(\ep)-sin(\al)*cos(\ep)*cos(\x)});
  \draw[sea!60,line width=0.4pt] plot[domain=-90:90,samples=60,smooth] ({cos(\ep)*sin(\x)},{-cos(\al)*sin(\ep)-sin(\al)*cos(\ep)*cos(\x)});
  \draw[seadeep,line width=0.8pt] plot[domain=-90:90,samples=60,smooth] ({sin(\x)},{-sin(\al)*cos(\x)});
  \draw[muted,dashed,line width=0.4pt] plot[domain=90:270,samples=60,smooth] ({sin(\x)},{-sin(\al)*cos(\x)});
  \draw[muted] (0,0) circle (1);
  \fill[seadeep] (0,{cos(\al)}) circle (1.2pt);
  \node[font=\small] at (0,-1.3) {band $\abs{\Phi}<\eps$};
\end{tikzpicture}\hspace{1.6cm}%
\begin{tikzpicture}[x=1.4cm,y=1.4cm,baseline=0pt]
  \def\al{18}\def\ep{14}
  \fill[seamist] plot[domain=-71.5:90,samples=60,smooth] ({cos(\x)*sin(\ep)},{cos(\al)*sin(\x)-sin(\al)*cos(\x)*cos(\ep)}) -- plot[domain=90:-71.5,samples=60,smooth] ({-cos(\x)*sin(\ep)},{cos(\al)*sin(\x)-sin(\al)*cos(\x)*cos(\ep)}) -- cycle;
  \draw[sea!60,line width=0.4pt] plot[domain=-71.5:90,samples=60,smooth] ({cos(\x)*sin(\ep)},{cos(\al)*sin(\x)-sin(\al)*cos(\x)*cos(\ep)});
  \draw[sea!60,line width=0.4pt] plot[domain=-71.5:90,samples=60,smooth] ({-cos(\x)*sin(\ep)},{cos(\al)*sin(\x)-sin(\al)*cos(\x)*cos(\ep)});
  \draw[seadeep,line width=0.8pt] plot[domain=-72:90,samples=60,smooth] ({0},{cos(\al)*sin(\x)-sin(\al)*cos(\x)});
  \draw[muted,line width=0.4pt] plot[domain=-90:90,samples=60,smooth] ({sin(\x)},{-sin(\al)*cos(\x)});
  \draw[muted,dashed,line width=0.4pt] plot[domain=90:270,samples=60,smooth] ({sin(\x)},{-sin(\al)*cos(\x)});
  \draw[muted] (0,0) circle (1);
  \fill[seadeep] (0,{cos(\al)}) circle (1.2pt);
  \node[font=\small] at (0,-1.3) {lune $\abs{\Lambda-\lambda_0}<\eps$};
\end{tikzpicture}
\caption{Two ways to close in on a great circle. The band around the equator has the same width everywhere, so given the band the longitude is uniform. The lune around a meridian is widest at the equator, so given the lune the latitude has density $\frac12\cos\varphi$.}
\label{fig:kolmogorov:sphere}
\end{figure}

Bertrand's paradox dissolves in the same way, by computing the three measures that the three recipes declare.

\begin{proposition}[Three measures on chords]\label[proposition]{prop:kolmogorov:chords}
Let $D$ be the distance from the center of the unit circle to a random chord. Under the recipes (i), (ii), and (iii) of \cref{sec:kolmogorov:chords}, $D$ has the densities $\frac2\pi(1-t^2)^{-1/2}$, $1$, and $2t$ on $[0,1)$, and $\kolP(D<\frac12)$ equals $\frac13$, $\frac12$, and $\frac14$.
\end{proposition}

\begin{proof}
Recipe (ii) makes $D$ uniform. Under recipe (iii) the midpoint $M$ is uniform in the disk and $D=\abs{M}$, so $\kolP(D\le t)=t^2$. Under recipe (i), put the first endpoint at angle $0$ and the second at angle $\Psi$, uniform on $[0,2\pi)$; then $D=\abs{\cos(\Psi/2)}$, and since $\Psi/2$ is uniform on $[0,\pi)$, $\kolP(D\le t)=1-\frac2\pi\arccos t$. The degenerate cases, $\Psi=0$ in recipe (i) and $M$ at the center in recipe (iii), have probability zero. Differentiating gives the densities, and $t=\frac12$ gives the probabilities, since $\arccos\frac12=\frac\pi3$.
\end{proof}

The three answers are probabilities of one event, $\{D<\frac12\}$, under three probability measures on the set of chords. The axioms do not choose among them; they make the choice a declared parameter. \textcite{Jaynes1973} proposed a way to choose: require the answer to be unchanged by rotations, dilations, and translations of the circle relative to the random lines, and only recipe (ii) survives.\sidenote{Jaynes also reported a trial in which straws were tossed onto a circle; the observed chord lengths agreed with recipe (ii).} That is a different move, a symmetry imposed on the model, which the axioms make expressible but do not supply; compare \cref{ch:conformal}.

\section{Costs and limits}\label{sec:kolmogorov:costs}

\paragraph{Where did the difficulty go?} Into measure theory, which \cref{ch:lebesgue} paid for, and into three existence theorems, each proved once. Carathéodory's extension builds a probability from its values on simple events; it is how Lebesgue measure becomes the law of fair coin tosses. Kolmogorov's extension theorem builds infinite families of random variables from consistent finite-dimensional laws. The Radon--Nikodym theorem, or the projection theorem with completeness, builds conditional expectations. After these payments the theorems of this chapter are short: the strong law took a fourth-moment computation, the zero--one law an approximation, and the tower property one line. Before them, each model was encoded by hand into lengths of intervals, as in the aside, and each theorem was proved inside one encoding. Paying once for every model is the bargain.

\paragraph{What did not dissolve.} First, modeling. The axioms say what a probability is, not which one describes a given random phenomenon; Bertrand's question is still unanswerable as posed, and symmetry, data, or a decision must supply the measure. Second, interpretation. Whether a probability is a limiting frequency, as \textcite{vonMises1919} held with his collectives, or a degree of belief, as de~Finetti held, the axioms do not say; de~Finetti also rejected countable additivity and built probability on finitely additive measures \parencite{deFinetti1974}. Third, conditioning on null events.\index{regular conditional probability} A conditional expectation is defined only almost surely. A \term{regular conditional distribution} of a random variable $X$ given $\cG$, a probability measure on the values of $X$ for each $\omega$ that represents $\kolP(X\in\cdot\mid\cG)$, exists when $X$ takes values in a Polish space \parencite[§4.1]{Durrett2019}; but regular conditional probabilities on an arbitrary probability space need not exist, and \textcite{Dieudonne1948} gave the first example. Even a regular version takes values on a given null set that depend on the $\sigma$-algebra, as the sphere shows. Fourth, uncountable families. In $\R^{[0,1]}$ with the $\sigma$-algebra of \cref{thm:kolmogorov:extension}, every event depends on countably many coordinates, so the set of continuous paths is not an event; Brownian motion with continuous paths needs a further construction, which \textcite{Wiener1923} had given ten years before the \emph{Grundbegriffe}.

\paragraph{Other seas.} Game-theoretic probability \parencite{ShaferVovk2019} takes betting rather than measure as primitive, following Ville's work on martingales: the strong law becomes the statement that a gambler has a strategy which never risks bankruptcy and makes him infinitely rich unless the average of the outcomes minus the announced probabilities tends to zero. Algorithmic randomness, starting with \textcite{MartinLof1966}, defines individual random sequences, the objects von Mises wanted, as those that pass every effective statistical test. Quantum probability replaces commuting random variables by operators, and the $\sigma$-algebra of events by a lattice of projections.

\section{The move, distilled}\label{sec:kolmogorov:distilled}

\begin{move}{Axiomatization}
\movepart{In one sentence} When an intuitive notion keeps producing contradictory answers, make the structure that every correct computation uses into the definition; the contradictions become different values of a parameter that the definition forces you to declare.
\movepart{Triggers} Two reasonable computations that give two answers; arguments ``by symmetry'' or ``at random'' that never say with respect to what; theorems proved anew for each model; disputes that never end because the parties answer different questions.
\movepart{Recipe} Collect the computations that work. Extract the structure they use (here, a $\sigma$-algebra of events, a countably additive measure of total mass one, independence as a product rule, and conditioning as projection). Make that structure the definition. Re-derive each paradox inside it and name the parameter it left implicit. Prove the basic theorems once, for every model.
\movepart{Anti-pattern} An axiom that settles a paradox by decree, declaring for instance that a random chord has a uniformly distributed distance, instead of exposing the choice; or axioms so weak that the theorems one wants cannot be proved, or so strong that the intended examples are excluded.
\end{move}

\dissolutiontest{Borel's normal number theorem, for every base and block, follows in a few lines from one strong law instead of a counting-and-covering argument per case; the tower property takes one line instead of a computation with densities.}{Conditioning needs no densities, continuity, or positivity; no model has to be encoded by hand into the unit interval.}{Coin tossing, dice, continuous measurements, and infinite sequences are one kind of object; expectation is Lebesgue's integral; conditional expectation is orthogonal projection.}{Kolmogorov's zero--one law; conditioning on an arbitrary $\sigma$-algebra, the starting point of martingale theory; a resolution of the sphere paradox.}

Axiomatization is the second move of Part~I. \Cref{ch:matroid} makes it for greedy algorithms: the exchange axiom is the structure that every greedy proof uses, and once it is the definition, one theorem covers spanning trees, linear independence, and scheduling. Both chapters pay in the same currency, by checking the axioms for each new example; here that means checking measurability and constructing the space. The move also leans on \cref{ch:lebesgue}: without the completed integral there would be no measure to axiomatize, and the strong law would have no almost-everywhere limits to speak of.

Probability spaces are the setting of \cref{ch:vc,ch:lecam,ch:conformal}. Exchangeability, the free symmetry of \cref{ch:conformal}, is itself a modeling assumption of exactly the kind the axioms force one to declare, and the hidden-parameter diagnosis returns in the research exercise below, where the parameter is a metric.

\begin{exercises}
\item\label{ex:kolmogorov:bertrand} \consolidation (Bertrand's measures.) (a) Show that the expected length of the chord is $\frac4\pi$, $\frac\pi2$, and $\frac43$ under recipes (i), (ii), and (iii). (b) Describe the line $\{x\cos\theta+y\sin\theta=p\}$ by the pair $(\theta,p)\in[0,\pi)\times\R$. Show that the measure $\dd\theta\dd p$ is invariant under rotations and translations of the plane, and that, restricted to the lines that meet the unit circle, it gives recipe (ii). (c) Recipe (iii) corresponds to the measure $\abs{p}\dd\theta\dd p$. Show that it is not invariant under translations.

\item\label{ex:kolmogorov:lunes} \consolidation (Bands and lunes.) Let $U$ be uniform on the sphere. (a) Show that, given the event $\{\abs\Phi<\eps\}$, the longitude is uniform, and that, given $\{\abs{\Lambda-\lambda_0}<\eps\}$ (the difference taken modulo $2\pi$), the latitude has density $\frac12\cos\varphi$, for every $\eps\in(0,\frac\pi2)$. (b) Let $N_\eps$ be the set of points at angular distance less than $\eps$ from the great circle that contains the meridian $\{\Lambda=\lambda_0\}$. Show that, given $\{U\in N_\eps\}$, the law of the latitude tends to the uniform law on $[-\frac\pi2,\frac\pi2]$ as $\eps\to0$. Which $\sigma$-algebra does this limit correspond to?

\item\label{ex:kolmogorov:slln} \consolidation (The strong law, quantitatively.) (a) In the setting of \cref{thm:kolmogorov:slln}, show that $\kolP(\abs{\frac1nS_n-\mu}\ge\eps)\le3K\eps^{-4}n^{-2}$, and deduce the weak law. (b) Let $X_1,X_2,\dots$ be independent with common distribution function $F$. Show that, almost surely, $\frac1n\#\{i\le n:X_i\le t\}\to F(t)$ for all rational $t$ simultaneously. (c) Show by an example that independence cannot simply be dropped.

\item\label{ex:kolmogorov:tail} \consolidation (Tail events.) Let $X_1,X_2,\dots$ be independent fair coin tosses with values $0$ and $1$, and $S_n=X_1+\dots+X_n$. (a) Which of these events are tail events: $\{X_n=1$ for infinitely many $n\}$, $\{\frac1nS_n\to\frac12\}$, $\{X_1=1\}$, $\{\limsup_n(S_n-\frac n2)/\sqrt n\ge1\}$? (b) Compute the probability of the first event directly, by continuity of $\kolP$. (c) Show that every random variable that is measurable with respect to the tail $\sigma$-algebra is almost surely constant.

\item\label{ex:kolmogorov:condexp} \consolidation (Working with conditional expectation.) Let $X\in L^2$ and let $\cG\subseteq\cF$ be a $\sigma$-algebra. (a) Show that if $X$ is independent of $\cG$, then $\E[X\mid\cG]=\E X$. (b) Show that $\E[X\mid\cG]$ minimizes $\E(X-W)^2$ over $\cG$-measurable $W\in L^2$, and, with $\Var(X\mid\cG)=\E[X^2\mid\cG]-\E[X\mid\cG]^2$, that $\Var X=\E\Var(X\mid\cG)+\Var\E[X\mid\cG]$. (c) Show that if $X'\in L^2$ and $X\le X'$ almost surely, then $\E[X\mid\cG]\le\E[X'\mid\cG]$ almost surely.

\item\label{ex:kolmogorov:entropy} \transfer\hard (Entropy from axioms.)\index{entropy} Let $H$ assign a number to every probability vector $(p_1,\dots,p_n)$, $n\ge1$, and suppose that (1) $H$ is symmetric in its arguments; (2) $H(p_1,\dots,p_n)=H(p_1+p_2,p_3,\dots,p_n)+(p_1+p_2)\,H\bigl(\frac{p_1}{p_1+p_2},\frac{p_2}{p_1+p_2}\bigr)$ whenever $p_1+p_2>0$; (3) $A(n)=H(\frac1n,\dots,\frac1n)$ is nondecreasing in $n$. (a) Show that $A(mn)=A(m)+A(n)$. (b) Show that $A(n)=c\log n$ for a constant $c\ge0$. (c) Deduce that $H(p)=-c\sum_ip_i\log p_i$ for every probability vector with positive rational entries. (d) Which parameter does $c$ represent? The argument is Shannon's \parencite{Shannon1948}; \textcite{Faddeev1956} showed that (3) can be replaced by a continuity assumption. \hint{for (a), first extend (2) to grouping into blocks; for (b), compare $s^k\le t^m<s^{k+1}$.}

\item\label{ex:kolmogorov:privacy} \transfer (Privacy as a definition.) \textcite{DinurNissim2003} showed that answering $O(n\log^2n)$ random subset-sum queries about $n$ secret bits with error $o(\sqrt n)$ lets a polynomial-time adversary reconstruct all but $o(n)$ of the bits. \textcite{DworkMcSherryNissimSmith2006} proposed a definition instead of a list of attacks: a randomized algorithm $M$ is \term{$\eps$-differentially private}\index{differential privacy} if $\kolP(M(D)\in S)\le e^\eps\kolP(M(D')\in S)$ for every event $S$ and all datasets $D,D'$ that differ in one record. (a) Let $\Delta>0$ and $\eps>0$, and let $q$ be a real-valued query with $\abs{q(D)-q(D')}\le\Delta$ for such pairs. Show that $q(D)+Z$, with $Z$ of density $\frac{\eps}{2\Delta}e^{-\eps\abs z/\Delta}$, is $\eps$-differentially private. (b) Show that running an $\eps_1$- and an $\eps_2$-differentially private algorithm with independent randomness is $(\eps_1+\eps_2)$-differentially private. (c) Show that $g\circ M$ is $\eps$-differentially private for every measurable $g$. (d) Which parameters does the definition force one to declare, and which question about the adversary does it make unnecessary? \hint{for (b), apply Fubini's theorem to sections, and note that $\mu(A)\le e^\eps\nu(A)$ for every $A$ implies $\int f\dd\mu\le e^\eps\int f\dd\nu$ for $f\ge0$.}

\item\label{ex:kolmogorov:random} \transfer (What is a random sequence?) Richard \textcite{vonMises1919} defined probability through \term{collectives}:\index{collective} infinite $0$--$1$ sequences whose limiting frequency of ones exists and is unchanged along every subsequence picked by an admissible selection rule. (a) Show that if every subsequence were admissible, then no sequence with limiting frequency strictly between $0$ and $1$ would be a collective. (b) Show that, for every countable family of infinite sets of indices fixed in advance, almost every sequence of fair coin tosses has limiting frequency $\frac12$ along each of them. (c) \textcite{MartinLof1966} called a sequence random if it lies outside every effectively null set: every set that can be covered, for each $k$, by a union of intervals with dyadic endpoints and total length at most $2^{-k}$, listed by a single algorithm that receives $k$ as input. Explain which parameter this definition declares that von Mises left open, and show that the random sequences in this sense have probability one. \hint{there are only countably many algorithms.}

\item\label{ex:kolmogorov:flat} \research\harder (Flat minima.) \textcite{DinhPascanuBengioBengio2017} observed that the flatness of a minimum is not a property of the function that a network computes. (a) For data $(x,y)$ with $\E x^2,\E y^2<\infty$ and $\E(x^+)^2>0$, let $L(w_1,w_2)=\E\bigl(y-w_2\max(w_1x,0)\bigr)^2$ for $w_1>0$. Show that $L(w_1,w_2)=g(w_1w_2)$ for a quadratic $g$, and that at a minimizer with $u^*=w_1w_2\ne0$ the Hessian of $L$ has the eigenvalues $0$ and $g''(u^*)(w_1^2+w_2^2)$. Conclude that along the curve $(\alpha w_1,w_2/\alpha)$, $\alpha>0$, which computes the same function, the largest eigenvalue takes every value in $[2g''(u^*)\abs{u^*},\infty)$. (b) Open: propose an axiomatization of flatness, a functional of the loss landscape that is invariant under the symmetries of the parametrization and still predicts generalization. Which hidden parameter, such as a metric or a prior on the parameters, must such a definition declare, and is there a choice for which flatness provably controls the generalization gap?
\end{exercises}

\notesandrefs

The \emph{Grundbegriffe} \parencite{Kolmogorov1933} is short and readable, and it is available in English \parencite{Kolmogorov1956}; Chapter~V, §2 contains the explanation of the sphere paradox. \textcite{ShaferVovk2006} trace its sources: Borel's denumerable probabilities, Cantelli's strong law, von Mises's collectives, Steinhaus's measure-theoretic model of coin tossing, product measures in the work of Łomnicki and Ulam, and a 1919 paper of Daniell that already contains the essential content of the extension theorem. Kolmogorov's own strong law, under the condition $\sum_n\Var X_n/n^2<\infty$, is in \textcite{Kolmogorov1930}.

\textcite{Williams1991} develops conditional expectation by projection, as in \cref{sec:kolmogorov:projection}, and proves the fourth-moment strong law; \textcite{Durrett2019} covers the extension theorem and regular conditional distributions. On Bertrand's chords see \textcite{Jaynes1973}. Game-theoretic probability is developed in \textcite{ShaferVovk2019}, and \textcite{Ville1939} is the critique of collectives from which it grew. The transfer exercises draw on \textcite{Shannon1948}, \textcite{Faddeev1956}, \textcite{DinurNissim2003}, and \textcite{DworkMcSherryNissimSmith2006}.

\printchapterbib
\end{refsection}
